% Separate extension of finite-puncture descent; no assertion of general finite generation.
\begin{corollary}[Proved: nonnormal schemes with controlled exceptional centres]
\label{D:nonnormal-controlled-boundary-quantitative-transfer}
Let $S$ be a proper scheme over an algebraically closed field $k$
of characteristic $p>0$, and let $T\subset S$ be closed, of
codimension at least two in every irreducible component which it
meets.  Let $U$ be a connected component of $S\setminus T$, and put
$Q=\pi_1(U)$.  Write $S_i$ for the normalizations of the irreducible
components of $S_{\mathrm{red}}$ which meet $U$, and write $T_i$
for the inverse image of $T$ in $S_i$.

For each $S_i$ of dimension at least two, assume that there is a
proper birational morphism $q_i:Y_i\to S_i$, with $Y_i$ normal and
projective, such that every prime divisor on $Y_i$ whose image is
contained in $T_i$ and has codimension at least two in $S_i$ has
zero-dimensional image.  Then all three quantitative conclusions
of Corollary~\ref{D:nonnormal-finite-puncture-quantitative-transfer}
hold for $Q$, with constants depending on this fixed pair and these
modifications.  Explicitly:
\begin{enumerate}
\item There is a constant $c$ such that
\[
             \dim_\Lambda H^1(Q,M)\le c\dim_\Lambda M
\]
for every finite field $\Lambda$ and every finite-dimensional
continuous $\Lambda[Q]$-module $M$.  For an open subgroup $J$ the
constant can be taken to be $c[Q:J]$.  The maximal prosolvable
quotient of $Q$ is topologically finitely generated.
\item There is a constant $C$ such that, for $q=p^f$ and $r\ge1$,
\[
 \#\bigl(\operatorname{Hom}_{\mathrm{cts}}
       (Q,\operatorname{GL}_r(\mathbf F_q))
       /\operatorname{GL}_r(\mathbf F_q)\bigr)
                 \le q^{Cr^2}.
\]
\item Fix $B\ge1$, and put $q=p^f$ with $f\ge1$.  Suppose
$\widetilde G\subset\operatorname{GL}_r(\mathbf F_q)$ has scalar
subgroup of order at most $B$, and let
$G\subset\operatorname{PGL}_r(\mathbf F_q)$ be its projective
image.  The number of homomorphisms $Q\to G$ up to conjugacy in
the ambient projective linear group is at most $q^{C_Br^2}$,
where $C_B$ is independent of $q,r,\widetilde G$.
For every fixed $\epsilon>0$, the quotient formed from all distinct
nonabelian simple quotients in any such family satisfying
$|G|\ge q^{\epsilon r^2}$ is topologically finitely generated.
\end{enumerate}
The omitted set is allowed to have positive dimension.  The
geometric condition on the normalized components is an explicit
hypothesis; it is not asserted for arbitrary projective
modifications.  No degree-two cohomology bound or full finite
generation of $Q$ is asserted.
\end{corollary}

\begin{proof}
Passing to the reduction leaves the finite etale category unchanged.
The normalization of $S_{\mathrm{red}}$ is finite because $S$ is
of finite type over a field.  Its inverse image $V$ of $U$ is a
finite surjective $U$-scheme.  Each nonempty open
$S_i\setminus T_i$ is integral and therefore connected, so it lies
over one connected component of $S\setminus T$.  Consequently the
connected components of $V$ are exactly the nonempty schemes
\[
                         V_i=S_i\setminus T_i
\]
which lie over $U$.

The finite map $S_i\to S$ preserves the dimensions of closed
irreducible subsets.  Thus $T_i$ has codimension at least two in
$S_i$; in particular $T_i$ is empty when $\dim S_i\le1$.
For every $i$ of dimension at least two,
Theorem~\ref{D:generic-surface-controlled-boundary} supplies a fixed
normal proper surface $X_i$ over an algebraically closed extension
of $k$, a finite set $F_i\subset X_i$, and a surjection
\[
                   \pi_1(X_i\setminus F_i)\twoheadrightarrow
                   \pi_1(V_i).
\]
The surface and extension field are independent of the covering
group and the coefficient field.  The established surface bounds
therefore give all three assertions for $\pi_1(V_i)$: use
degree-one inflation for cohomology and injective precomposition
for ordinary or ambient-projective representation classes.
Components of dimension at most one are proper and satisfy the
same bounds by the already established lower-dimensional cases.

For completeness we give the exact estimates which descend these
bounds.  Put $a=\#\pi_0(V\times_UV)$, a finite number.  If a finite
local system $M$ on $U$ has rank $r$, the kernel of
\[
 H^1(U,M)\longrightarrow\bigoplus_iH^1(V_i,M|_{V_i})
\]
is a quotient of a subspace of $H^0(V\times_UV,M)$.  Indeed an
$M$-torsor trivialized on $V$ is specified by its translation
cocycle on the double product, modulo changes of trivialization.
Finite-proper descent is effective and fully faithful for these
finite etale torsors \citep[Expose IX, Theorem 4.7]{C_SGA1}.
A section on a connected component is determined by one geometric
stalk, so this kernel has dimension at most $ar$.  Adding the
finitely many bounds on $V_i$ proves (1).  Shapiro's lemma gives
the stated open-subgroup bound, and
Proposition~\ref{D:cohomology-crown-criterion} gives its
prosolvable consequence.

For (2), put $P=\operatorname{GL}_r(\mathbf F_q)$.  Fix a tuple
of $P$-torsor classes on the components $V_i$ and representatives
for it.  On the double product the Isom scheme between the two
pullbacks is finite etale, with $|P|$ elements in each geometric
fibre.  There are therefore at most $|P|^a$ possible descent
isomorphisms.  The cocycle conditions and passage to isomorphism
classes can only decrease this number.  The componentwise bounds
and $|P|\le q^{r^2}$ yield the required estimate.

For (3), repeat this argument with
$P=\operatorname{PGL}_r(\mathbf F_q)$, restricting the possible
componentwise torsor classes to ones coming from $G$-valued
homomorphisms.  Extension of structure group identifies equality
of the resulting $P$-torsor classes precisely with ambient
$P$-conjugacy.  This is an equivalence relation on the original
$G$-valued homomorphisms even when $P$ does not normalize $G$.
For a fixed tuple it is legitimate to count all $P$-descent data,
including data whose final monodromy does not lie in $G$; this
only enlarges the upper bound.  The same factor $|P|^a$ proves
the estimate.  Distinct epimorphism kernels are separated by
ambient conjugacy, so the order lower bound converts it to a
fixed polynomial bound in $|G|$ on the multiplicities of the
specified simple quotients.  The generating-tuple argument of
Proposition~\ref{D:bounded-scalar-projective-growth} completes the
proof.  All descent used here is in degrees zero and one.
\end{proof}
